\begin{center}
	\section*{Homework 01 - 2.2, 2.3}
	Due Feb 13 \\
	Uzair Hamed Mohammed
\end{center}

%12, 16, 22, 38
%10, 12, 14

\subsection*{2.2: Sample Spaces and the Algebra of Sets}

\begin{enumerate}
	\item[12.] Consider the experiment of choosing coefficients for the quadratic equation \( ax^2 + bx + c = 0 \). Characterize the values of \(a, b\), and \(c\) associated with the event A: Equation has complex roots.
	
		\underline{Sol}: 
		\[
			\begin{aligned}
				\textrm{Sample space } S: \lbrace a, b, c \in \mathbb{R} \rbrace \\
				\textrm{Complex roots occur when discriminant } \boxed{b^2 -4ac < 0}
			\end{aligned}
		\]
	
	\item[16.] Sketch the regions in the xy-plane corresponding to \(A \cup B\) and \(A \cap B\) if 
		\[
			A = \lbrace (x, y): 0 < x < 3, 0 < y < 3 \rbrace
		\]
		
		and
		
		\[
			B = \lbrace (x, y): 2 < x < 4, 2 < y < 4 \rbrace
		\]
		
		\underline{Sol}: \vspace{2.0in}
	
	\item[22.] Suppose that each of the twelve letters in the word 
		\[
			T \ E \ S \ S \ E \ L \ L \ A \ T \ I \ O \ N
		\]
		
		is written on a chip. Define the events F, R and C as follows:
			\begin{itemize}
				\item[F: ] letters in first half of alphabet
				\item[R: ] letters that are repeated
				\item[V: ] letters that are vowels
			\end{itemize}
		Which chips make up the following events? 
			\begin{enumerate}
				\item[a.] \(F \cap R \cap V\)
				\item[b.] \(F^C \cap R \cap V^C\)
				\item[c.] \(F \cap R^C \cap V\)
			\end{enumerate}

		\underline{Sol}:
		\[
			\begin{aligned}
				F &= \lbrace A, E, I, L \rbrace \\
				R &= \lbrace E, I, L, S, T \rbrace \\
				V &= \lbrace A, E, I, O \rbrace
			\end{aligned}
		\]
		\begin{enumerate}
			\item \(F \cap R \cap V = \boxed{\lbrace E, I \rbrace}\)
			
			\item \(F^C \cap R \cap V^C = \boxed{\lbrace S, T \rbrace}\)
			
			\item \(F \cap R^C \cap V = \boxed{\lbrace A \rbrace}\)
		\end{enumerate}
		
	\item[38.] Let A, B, C be any three events defined on a sample space S. Let \(N(A), N(B), N(C), N(A \cap B), N(A \cap C), N(B \cap C), N(A \cap B \cap C)\) denote the numbers of outcomes in all the different intersections in which A, B, and C are involved. Use a Venn diagram to suggest a formula for \(N(A \cup B \cup C)\).
	
		\underline{Sol}:
		\[
			N(A \cup B \cup C) = \boxed{N(A) + N(B) + N(C) - N(A \cap B) - N(A \cap C) - N(B \cap C) + N(A \cap B \cap C)}
		\]
\end{enumerate}

\subsection*{2.3: The Probability Function}

\begin{enumerate}
	\item[10.] An urn contains twenty-four chips, numbered 1 through 24. One is drawn at random. Let A be the event that the number is divisible by 2 and let B be the event that the number is divisible by 3. Find \(P(A \cup B)\).
	
		\underline{Sol}:
		\[
			\begin{aligned}
				P(A \cup B) &= P(A) + P(B) - P(A \cap B) \\
				&= \frac{12}{24} + \frac{8}{24} - \frac{4}{24} \\
				&= \frac{16}{24} \\
				&= \boxed{\frac{2}{3}}
			\end{aligned}
		\]
	
	\item[12.] Events \(A_1\) and \(A_2\) are such that \(A_1 \cup A_2 = S\) and \(A_1 \cap A_2 = \emptyset\). Find \(p_2\) if \(P(A_1) = p_1, P(A_2) = p_2\), and \(3p_1 - p_2 = \frac{1}{2}\).
	
		\underline{Sol}:
		\[
			\boxed{p_1 = \frac{3}{8}, p_2 = \frac{5}{8}}
		\]
	
	\item[14.] Three events- A, B, C- are defined on a sample space S. Given that \(P(A) = 0.2, P(B) = 0.1 \), and \(P(C) = 0.3\), what is the smallest possible value for \(P \lbrack (A \cup B \cup C)^C \rbrack \)?
	
		\underline{Sol}:
		\[
			\begin{aligned}
				A \cap B &= A \cap C = B \cap C = \emptyset \\
				&= 1 - 0.2 - 0.1 - 0.3 \\
				&= \boxed{0.5}
			\end{aligned}
		\]
\end{enumerate}